\chapter{Research Methodology}
\label{ch:methodology}

\section{Research strategy}

This study uses a design-science and software-engineering strategy. The artefact is a Godot 4.6 plugin and executable demo; knowledge is produced by connecting design requirements from behaviour-tree semantics and visualisation literature to an implemented system, then measuring correctness, display properties, authoring interaction and developer-task usability. Runtime profiling is retained as secondary engineering evidence rather than the main criterion of success. Iterative development was appropriate because editor interactions exposed faults that were not visible in isolated algorithms, including coordinate-space errors, stale transforms, duplicated edges and partially written live-debug snapshots.

The process had five stages. First, supervisor requirements and the dissertation guidance were converted into acceptance criteria. Second, a minimal executable behaviour-tree core and game integration were established. Third, editor operations and live debugging were added. Fourth, visual techniques were implemented behind independent switches, with a disable/reset path tested after each addition. Fifth, deterministic workloads and layered tests were used to compare the completed artefact against its own baseline. Later tasks were followed by full regression because a change to graph rebuilding or resource copying could invalidate an earlier feature.

\section{Experimental factors and datasets}

Three deterministic generated trees contain 31, 121 and 364 nodes. They represent a small tree that is comfortably editable, a medium tree exceeding the approximately 50-node readability concern raised during supervision, and a large stress case. Each tree is generated by the same structural rule so that scale, rather than unrelated content, is the main difference. A separate ecological fixture contains 241 resource nodes: 202 canvas cards and 39 attached decorators. Unlike the generated trees, it is the resource used by the playable enemy and includes recovery, directional melee, obstacle jumping, ladder climbing, ranged projectiles, pressure, chase, last-known-position search, return, patrol and idle behaviour.

The display experiment uses the conventional full-detail node--link graph as the baseline. The principal conditions are compact cards, semantic zoom, search highlight, subtree focus, subtree collapse, fisheye and enhanced minimap. Additional switch-cycle checks cover all 19 features. Measurements are defined as follows:

\begin{itemize}
    \item \emph{visible ratio}: visible nodes divided by total nodes;
    \item \emph{card area}: sum of visible card minimum-width multiplied by minimum-height;
    \item \emph{information fields}: representative fields retained at the active semantic-detail level;
    \item \emph{dimmed nodes}: nodes whose salience is reduced by search or runtime-path focus;
    \item \emph{focus scale gain}: internal magnification of the fisheye target;
    \item \emph{overview coverage}: nodes represented in the minimap divided by total nodes;
    \item \emph{operation time}: elapsed time for the measured rebuild, update or switch cycle on the test machine.
\end{itemize}

These are instrumented interface measures, not psychological measures. They identify how the implementation changes information density and context. The human study instead records task completion time, incorrect selections, zoom and pan counts, task success, and seven-point readability and cognitive-load ratings.

\section{Playable-tree display experiment}

The ecological display experiment loads the fixed 241-node playable resource in a 1600 by 900 pixel \texttt{SubViewport} using Godot 4.6's OpenGL Compatibility renderer on an NVIDIA GeForce RTX 5070 Laptop \ac{GPU}. The six conditions are Baseline, Compact Cards, Optimized Overview, Optimized Search, Subtree Focus and Context Collapse. Search uses \texttt{Patrol Route Choice}; focus uses \texttt{11 Layered Patrol}. Each condition receives three warm-ups and 30 recorded repetitions. Condition order rotates once per repetition, so every condition occupies each of the six sequence positions five times, reducing time and thermal drift without introducing an unreproducible random order.

Every observation begins from the same expanded tree, empty query, 100\% zoom and no focus. The script records resource-node and card counts, visible ratio, graph-bounds area, card area, overlap pairs, information fields, dimmed cards and interaction time. Raw data contain 180 observations; the summary reports median, first and third quartiles, and 95th percentile. Because the conditions trigger different operations, their time values describe response ranges rather than a common performance ranking. Assertions require exactly 241 resources and 202 baseline cards, zero overlaps, no saved-position mutation, at least 40\% compact-area reduction, at least 70\% overview-area reduction, at least 95\% non-target dimming, and at least 70\% card reduction for focus and collapse.

A second same-task comparison measures 240 pointer samples while dragging one card through a fixed path. The baseline and optimized versions use the same tree, card, path and renderer. One warm-up precedes five trials; median total and per-sample time are retained. This measures latency in an authoring action, not behaviour-tree runtime efficiency or human-perceived smoothness.

Runtime profiling crosses three tree sizes with 1, 10 and 50 runner instances and cache enabled/disabled. Every combination is warmed up and measured seven times. Tick budgets are adjusted so each condition performs a comparable bounded workload. The median mean-tick time is used because it is less sensitive than the arithmetic mean to transient scheduling delays. Correctness assertions require cached and uncached conditions to return the same terminal status.

\section{Software correctness and integration evaluation}

Testing is divided into five core layers. Resource/runtime tests validate structure, node semantics, decorators, actor calls, blackboard schemas, debug snapshots and cache invalidation. Editor GUI tests instantiate the plugin interface and exercise creation, dragging, connections, typed fields, save/reload, undo/redo and feature switches. Basic game integration tests verify player and enemy interaction. Complex-arena tests verify detection, reactive defence, directional melee, ranged projectiles, obstacle sensing and jumping, elevated-player ladder pursuit, loss of target, retreat and healing. Real-\ac{GPU} regression tests capture 1600 by 900 pixel views and assert visible graph properties that cannot be read from Godot's dummy renderer.

The pass policy is deliberately strict. A non-zero process exit, \texttt{FAIL:}, \texttt{ERROR:}, \texttt{SCRIPT ERROR:}, leaked object, resources left in use, access violation or native crash invalidates a suite. Warnings are examined by meaning: a deliberately missing actor method is a safety test and must return \textsc{failure}, while an editor file-scan warning after \texttt{--quit-after} is expected shutdown behaviour.

\section{Visual validation}

Automated assertions cannot detect every presentation fault. The visual protocol therefore renders stable scenes on an NVIDIA GeForce RTX 5070 Laptop \ac{GPU} using Godot's OpenGL 3.3 Compatibility renderer. Core screenshots cover baseline, type encoding, grayscale, colourblind palette, single connections, compact cards, fisheye enabled/disabled, collapse, search, runtime failure, live blackboard, schema editor, schema-key picker, orthogonal edges, bundled edges and the complex tree. Four additional controlled-study images cover the 121-node practice tree, the playable 241-node overview, a balanced search target and a deterministic seven-node gameplay path.

Visual inspection checks card overlap, port position, edge duplication, text clipping, stale scaling, viewport centring, active-path continuity, dimming reset and minimap coverage. This was essential in diagnosing a fisheye implementation that appeared correct by state variables but used an off-centre transform and interacted incorrectly with wheel zoom.

\section{Human-study design}

The pre-registered within-participant study compares Baseline, Compact, Collapse, Fisheye, Search and Minimap on the playable 241-node tree. Twelve participants complete three game-development task types under all six conditions: locating an action from a gameplay requirement, reporting a seven-node \textsc{running} path, and increasing an attached decorator's duration by 0.10 seconds. This produces 18 formal trials per participant and 216 total. A separate 121-node tree provides one unscored practice trial per task.

Six matched action/path targets cover obstacle traversal, vertical pursuit, ranged suppression, mid-range pressure, last-known-position search and patrol. Six edit targets are real decorators attached to actions. Method order uses six balanced sequences, each assigned twice; task-block order uses three rotations, each assigned four times. Target index is \((m+p)\bmod 6\), where \(m\) is method index and \(p\) is participant index. Thus each participant sees six targets per task and every method--target combination occurs exactly twice across the sample, limiting learning from repeatedly finding one card.

The primary outcomes are 180-second-capped completion time and success. Secondary outcomes are incorrect selections, zooms, pans, total navigation, and seven-point readability and cognitive-load ratings. Each task is analysed separately. Continuous and ordinal outcomes use Friedman tests followed, only after a significant omnibus result, by exact paired Wilcoxon comparisons against Baseline with Holm correction. Success uses Cochran's Q followed by corrected exact McNemar comparisons. Median [Q1, Q3], paired effect sizes, bootstrap intervals, exclusions, failures and timeouts are retained. Fewer than eight complete participant sets restricts reporting to descriptives.

Participants use anonymous identifiers, provide consent and may withdraw. No names or contact details are stored. Recruitment must follow University ethics and supervisor guidance. At the time of writing, the generated workbook contains 216 planned rows and zero participant observations; therefore no inferential human-performance result appears in this draft.

\section{Validity and reproducibility}

Internal validity is strengthened by fixed resources, deterministic targets, independent switches, 30 repeated ecological-display samples, balanced condition/target orders and full regression after integration. Construct validity remains limited because visible area and field count are measures of density rather than readability. External validity is limited to one Windows laptop, Godot 4.6, the fixed generated fixtures and one implemented 2D game. Timing values are therefore reported with hardware, renderer, viewport, warm-up and workload rather than presented as universal constants.

Reproducibility materials include raw CSV files, screenshots, test scripts, formula workbooks, a clean-install packaging script and an MIT-licensed versioned plugin archive. Human-study materials are provided separately from automated evidence so that uncollected participant data cannot be mistaken for results.
